% ============================================================
% Tabelle 6.1: Architekturmechanismus -> MAF-Grundlage -> eigene
% Ergänzung -> Vertiefung — v01 (Claude, 03.09.)
% In Overleaf ablegen als: tables/maf-zuordnung.tex
% Quelle: probeschnitt-vertrag-6-1-6-2.md §2 (v02-Endform, FROZEN;
% Zeilen 8+9 = Kollegen-Korrektur „Anknüpfungspunkte statt —").
% v02 (03.09., 6.1-Kollegen-Runde, alle 3 Punkte gegengeprüft):
% ① NEUE Zeile Evidenz-/Provenienzmodell (5.3→6.2) — 2.4.6 führt
%   Provenienz als eine der ZWEI größten Framework-Lücken
%   (chap2.4.6.tex:78–83; Tabelle 2.5 Zeile wörtlich übernommen aus
%   chap2.4.6table.tex:75–78); fehlte in v01 (A3-Loch).
% ② Steward-Zeile: Framework-/Eigenanteil sauber getrennt (V1-Fund:
%   MAF = SerializeSession + IChatReducer; WIR = Ablage
%   state/steward/sessions/ + Reducer-Wahl; StewardChatRunner.cs).
% ③ „Lauf-Anker" gestrichen (undefinierter Begriff; Detail → 6.5).
% Damit 10 Zeilen. Caption an die neue Prüf-Formulierung angeglichen.
% v03 (03.09., Gesamtpass-Kollegen-Runde): Steward-Zeile
% „governancegeschützten Abläufen" -> „kontrolliert bereitgestellten
% Fähigkeiten" — Nachzug der 6.6-v02-Korrektur (Lese-Tools sind
% nicht gate-geschützt); Tabelle war stale geblieben.
% Zahlenfrei (V2-Regel: Nutzungszahlen erst am Freeze-Commit).
% Abschnitts-Spalte Kap.-5-seitig als Klartext „5.x" (wie die
% Kap.-4/5-Praxis; \ref-Pass = P2-Bündel).
% ============================================================

\begin{table}[!htb]
  \centering
  \small
  \begin{tabular}{p{0.26\textwidth}p{0.27\textwidth}p{0.27\textwidth}p{0.08\textwidth}}
    \hline
    \textbf{Architekturmechanismus (Kap.~\ref{chap:loesungsarchitektur})} & \textbf{MAF-Grundlage} &
    \textbf{eigene Ergänzung} & \textbf{Vert.} \\
    \hline
    Typisierte Workflow-Komposition (\ref{sec:logische-gesamtarchitektur}) &
    typisierte Nachrichten, bedingte Kanten,
    Verteilung und Zusammenführung &
    wiederverwendete Stufen-Executors und pfadspezifische Komposition & \ref{sec:realisierung-agent-executor-workflow} \\
    Verantwortungsverteilung agentisch/deterministisch (\ref{sec:logische-gesamtarchitektur}, konkretisiert \ref{sec:saeule-semantische-transformation}) &
    Agent-Abstraktion, Chat-Pipeline, Structured Output &
    Prompt- und Prüfdesign, formale Prüfregeln & \ref{sec:realisierung-agent-executor-workflow} \\
    Prüf-Reparatur-Schleife (\ref{sec:saeule-semantische-transformation}) &
    konditionale Kanten mit Rückführung &
    gemeinsame Entscheidungsregel und Versuchshistorie & \ref{sec:realisierung-agent-executor-workflow} \\
    Evidenz- und Provenienzmodell (\ref{sec:evidenzgebundene-verarbeitung}) &
    Structured Outputs und typisierte Nachrichten; kein fachliches
    Provenienzmodell &
    quellstandsbezogene Unit- und lauflokale Claim-Kennungen,
    Quellenbindung und Relationsregeln & \ref{sec:realisierung-agent-executor-workflow} \\
    Menschliche Entscheidungspunkte (\ref{sec:governance-mutation}) &
    Request Ports (Anfrage/Antwort) &
    Antwortregeln je Gate und Betriebsart,
    Review-Schicht und Einreichung & \ref{sec:realisierung-human-gates} \\
    Durable Ausführung (\ref{sec:durable-ausfuehrung-agentische-bedienung}) &
    Checkpoints, Wiederaufnahme, Sub-Workflow-Bindung &
    portfreie Teilworkflows mit Strukturprüfung & \ref{sec:realisierung-durabilitaet} \\
    Agentische Bedienebene (\ref{sec:durable-ausfuehrung-agentische-bedienung}) &
    Agent, Werkzeugregistrierung und Freigabefunktion,
    Sitzungsspeicherung und Verlaufsreduktion &
    Werkzeugraum aus bestehenden kontrolliert bereitgestellten
    Fähigkeiten, dateibasierte Sitzungsablage und gewählte
    Reduktionsstrategien &
    \ref{sec:realisierung-steward} \\
    Fachliche Zustandshaltung (\ref{sec:core-projektfortschreibung}) &
    Sitzungs-/Workflowzustand und Checkpoints; kein fachliches
    Zustandsmodell &
    Core-Repository, fachliches Zustandsmodell, Speichernaht,
    Invarianten, Historie & \ref{sec:realisierung-persistenz-aussenkopplung} \\
    Projektion und Rückweg (\ref{sec:projektionen-kontrollierte-aussenkopplung}) &
    Werkzeug- und Integrationsmechanismen; kein fachlicher
    Projektionsvertrag &
    deterministische Projektion und Sammlung; bedingt
    modellgestützte Aufbereitung im Rückweg; externe Freigabe & \ref{sec:realisierung-persistenz-aussenkopplung} \\
    Beobachtbarkeit (\ref{sec:evidenzgebundene-verarbeitung}) &
    Telemetrie-Anbindung der Chat-Pipeline &
    dateibasierte Ablage je Lauf und Metrikaggregation & \ref{sec:realisierung-observability} \\
    \hline
  \end{tabular}
  \caption[MAF-Grundlagen und projektspezifische Ergänzungen]{
    Architekturmechanismen, MAF-Grundlagen und projektspezifische
    Ergänzungen; die letzte Spalte nennt den Vertiefungsabschnitt.}
  \label{t:maf-zuordnung}
\end{table}
